\chapter{System Design}

\section{Proposed System Architecture}

GAIA is organised in four layers (Figure~\ref{fig:architecture}): the rover, the Flutter mobile application, the backend
services and the data layer. Two outside services are used: the Gemini API (chat, vision and embeddings) and web search.

\begin{figure}[H]
\centering
\scaledtikz{\begin{tikzpicture}
% Mobile application layer
\node[box,text width=21mm] (a1) at (-5.1,7.2) {Rover control,\\sensor switches};
\node[box,text width=21mm] (a2) at (-2.55,7.2) {Sensors\\view};
\node[box,text width=21mm] (a3) at (0,7.2) {Chat\\(streamed)};
\node[box,text width=21mm] (a4) at (2.55,7.2) {Diagnose};
\node[box,text width=21mm] (a5) at (5.1,7.2) {Live view};
\node[zone,fit=(a1)(a2)(a3)(a4)(a5)] (zapp) {};
\node[zlbl] at (zapp.north) {FLUTTER MOBILE APPLICATION};
% Backend layer
\node[box,text width=27mm] (b1) at (-4.8,3.3) {\textbf{Ingestion}\\validate, switches,\\windows, chunks};
\node[box,text width=27mm] (b2) at (-1.6,3.3) {\textbf{Pipeline A}\\planner, SQL facts,\\hybrid RAG, answer};
\node[box,text width=27mm] (b3) at (1.6,3.3) {\textbf{Pipeline B}\\VLM, research\\agent, writer};
\node[box,text width=27mm] (b4) at (4.8,3.3) {\textbf{Live activity}\\event log, SSE\\stream, state};
\node[zone,fit=(b1)(b2)(b3)(b4)] (zback) {};
\node[zlbl] at (zback.north) {BACKEND SERVICES (FASTAPI, PYTHON)};
% Data layer
\node[db,text width=24mm] (d1) at (-3.8,-0.9) {Sensor log\\readings, windows,\\switch events};
\node[db,text width=24mm] (d2) at (0,-0.9) {\textbf{Knowledge}\\\textbf{chunks}\\vector + text};
\node[db,text width=24mm] (d3) at (3.8,-0.9) {Diagnoses,\\chat memory,\\documents};
\node[zone,fit=(d1)(d2)(d3)] (zdata) {};
\node[zlbl] at (zdata.south) {POSTGRESQL 16 + PGVECTOR};
% Rover and external services
\node[gbox,text width=21mm] (rover) at (-9.2,3.3) {\textbf{Rover}\\ESP32-CAM,\\Nano, sensors};
\node[gbox,text width=25mm] (gem) at (9.6,4.4) {\textbf{Gemini API}\\chat, vision,\\embeddings};
\node[gbox,text width=25mm] (web) at (9.6,2.0) {\textbf{Web search}\\DuckDuckGo or\\Tavily};
\draw[arr] (zapp) -- node[lbl,right]{HTTP, SSE} (zback);
\draw[arr] (zback) -- (zdata);
\draw[arr] (rover) -- node[lbl,above]{telemetry} (zback.west);
\draw[barr] (rover.north) |- node[lbl,pos=0.75,above]{direct: control, video, stills} (a1.west);
\draw[arr] (zback.east |- gem) -- (gem.west);
\draw[arr] (zback.east |- web) -- (web.west);
\end{tikzpicture}}

\caption{Layered architecture of GAIA.}
\label{fig:architecture}
\end{figure}

\subsection{Design Principles}

\begin{enumerate}
\item \textbf{Numbers come from code, words come from the model.} Sensor statistics are computed by SQL and by deterministic
      templates. The language model receives these facts and phrases them. It is never asked to compute an average or a date.
\item \textbf{Two independent pipelines.} Pipeline B never reads sensor data or chat history, and Pipeline A never sees image pixels.
      The only connection is that Pipeline B \emph{writes} text chunks that Pipeline A may later \emph{read}. There is no control loop.
      A test records every SQL statement issued during a diagnosis and checks that none touches the sensor or chat tables.
\item \textbf{No feedback from chat into the store.} Chat answers are saved as conversation memory only. They are not embedded into the
      knowledge store, so a model's mistake cannot return later as ``farm knowledge''.
\item \textbf{Rover as a thin layer.} The rover samples, queues and posts. All interpretation happens in the backend.
\item \textbf{Degrade gracefully.} Without a model key, chat lists the retrieved facts and the offline hashing embedder keeps retrieval working.
\item \textbf{Explain gaps from records.} Which sensors were on, off or silent is recorded by code, so the assistant never has to guess.
\end{enumerate}

\subsection{Network Topologies}

The phone talks \emph{directly} to the rover for control, live video and still capture; none of this passes through the backend.
The rover pushes sensor readings to the backend by itself. The phone talks to the backend for sensor views, chat and
diagnosis (Figure~\ref{fig:deployment}). If the phone is connected to the rover's own hotspot, the rover has no route to the
backend. The app then offers a \emph{relay}: it reads the rover's latest reading every five seconds and forwards it to the same
backend endpoint, queueing up to one hour of readings while the backend is unreachable. The rover's firmware supports three
Wi-Fi modes (station with hotspot fallback, hotspot only, or hotspot plus station).

\begin{figure}[H]
\centering
\scaledtikz{\begin{tikzpicture}
\node[box,text width=38mm,minimum height=24mm] (phone) at (0,0) {\textbf{Smartphone}\\Android\\Flutter application (APK)\\built by GitHub Actions};
\node[box,text width=44mm,minimum height=30mm] (srv) at (9.4,0) {\textbf{Laptop or server}\\FastAPI (uvicorn, port 8000)\\PostgreSQL 16 + pgvector\\uploaded photos on disk\\background ingestion job};
\node[box,text width=44mm,minimum height=24mm] (rov) at (4.7,-6.0) {\textbf{Rover}\\ESP32-CAM: HTTP port 80,\\MJPEG stream port 81\\Arduino Nano, sensors};
\node[dsh,text width=32mm] (rt) at (4.7,3.2) {\textbf{Wi-Fi router}\\or rover hotspot};
\node[gbox,text width=36mm] (gem) at (16.0,1.6) {\textbf{Gemini API}\\(HTTPS)};
\node[gbox,text width=36mm] (web) at (16.0,-1.6) {\textbf{Web search}\\DuckDuckGo or Tavily\\(HTTPS)};
\draw[barr] (phone) -- node[lbl,above]{HTTP :8000} (srv);
\draw[barr] (phone.south) |- node[lbl,pos=0.3,left]{HTTP :80, :81} (rov.west);
\draw[arr] (rov.east) -| node[lbl,pos=0.8,right]{POST /telemetry} (srv.south);
\draw[arr] (srv.east) ++(0,0.9) -- (gem.west);
\draw[arr] (srv.east) ++(0,-0.9) -- (web.west);
\draw[darr] (phone.north) |- (rt.west);
\draw[darr] (srv.north) |- (rt.east);
\end{tikzpicture}}

\caption{Deployment view: devices, ports and protocols.}
\label{fig:deployment}
\end{figure}

\section{System Block Diagram}

Figure~\ref{fig:blockdiagram} follows the data from the sensors to the knowledge store (row A), from a question to a streamed
answer (row B) and from a photograph to a stored diagnosis (row C). The knowledge store is the single meeting point.

\begin{figure}[H]
\centering
\scaledtikz{\begin{tikzpicture}
\node[zlbl,anchor=west] at (-1.2,1.3) {A. SENSING AND INGESTION};
\node[gbox,text width=23mm] (r1) at (0,0) {Sensors\\(DHT11, MQ,\\soil probe)};
\node[box,text width=23mm] (r2) at (3.5,0) {Nano: sample,\\apply switches};
\node[box,text width=23mm] (r3) at (7.0,0) {ESP32: queue,\\push every 15 s};
\node[box,text width=23mm] (r4) at (10.5,0) {POST /telemetry:\\validate, store};
\node[box,text width=23mm] (r5) at (14.0,0) {15-min windows,\\text chunks,\\embeddings};
\draw[arr] (r1) -- (r2); \draw[arr] (r2) -- node[lbl,above]{UART} (r3);
\draw[arr] (r3) -- node[lbl,above]{Wi-Fi} (r4); \draw[arr] (r4) -- (r5);
\node[zlbl,anchor=west] at (-1.2,-2.0) {B. QUESTION ANSWERING (PIPELINE A)};
\node[gbox,text width=23mm] (q1) at (0,-3.3) {Farmer's\\question};
\node[box,text width=23mm] (q2) at (3.5,-3.3) {App: POST\\/chat/stream};
\node[box,text width=23mm] (q3) at (7.0,-3.3) {Plan, SQL facts,\\hybrid retrieval};
\node[box,text width=23mm] (q4) at (10.5,-3.3) {Gemini answer,\\citation check};
\node[box,text width=23mm] (q5) at (14.0,-3.3) {SSE tokens\\to the app};
\draw[arr] (q1) -- (q2); \draw[arr] (q2) -- (q3); \draw[arr] (q3) -- (q4); \draw[arr] (q4) -- (q5);
\node[zlbl,anchor=west] at (-1.2,-5.3) {C. CROP DIAGNOSIS (PIPELINE B)};
\node[gbox,text width=23mm] (c1) at (0,-6.6) {Photo: rover\\/capture or\\phone camera};
\node[box,text width=23mm] (c2) at (3.5,-6.6) {App: POST\\/diagnoses};
\node[box,text width=23mm] (c3) at (7.0,-6.6) {VLM call:\\strict JSON};
\node[box,text width=23mm] (c4) at (10.5,-6.6) {Research agent:\\web search};
\node[box,text width=23mm] (c5) at (14.0,-6.6) {Write diagnosis\\and findings\\chunks};
\draw[arr] (c1) -- (c2); \draw[arr] (c2) -- (c3); \draw[arr] (c3) -- (c4); \draw[arr] (c4) -- (c5);
\node[db,minimum height=28mm,text width=20mm] (kb) at (18.2,-3.3) {\textbf{Knowledge}\\\textbf{store}\\pgvector +\\full-text};
\draw[arr] (r5.east) -| (kb.north);
\draw[arr] (c5.east) -| (kb.south);
\draw[arr] (kb.west) -- node[lbl,above]{read} (q5.east);
\end{tikzpicture}}

\caption{System block diagram with the three data paths.}
\label{fig:blockdiagram}
\end{figure}

\subsection{Use Cases}

\begin{figure}[H]
\centering
\scaledtikz{\begin{tikzpicture}
\tikzset{uc/.style={nohyph,draw,ellipse,font=\dgfont,align=center,text width=44mm,inner sep=1pt,minimum height=9mm,fill=white}}
\def\actor#1#2#3{\draw (#1,#2) circle (0.2);\draw (#1,#2-0.2) -- (#1,#2-0.95);\draw (#1-0.35,#2-0.5) -- (#1+0.35,#2-0.5);\draw (#1,#2-0.95) -- (#1-0.3,#2-1.55);\draw (#1,#2-0.95) -- (#1+0.3,#2-1.55);\node[font=\dgfont\bfseries,align=center] at (#1,#2-2.1) {#3};}
\draw[rounded corners=3pt] (2.3,1.0) rectangle (10.4,-14.2);
\node[zlbl] at (6.35,0.8) {GAIA SYSTEM};
\node[uc] (u1) at (6.35,0) {Drive the rover, view live video};
\node[uc] (u2) at (6.35,-1.25) {Switch sensors on or off};
\node[uc] (u3) at (6.35,-2.5) {View sensor readings and charts};
\node[uc] (u4) at (6.35,-3.75) {Watch the Live view};
\node[uc] (u5) at (6.35,-5.0) {Ask about my farm (RAG chat)};
\node[uc] (u6) at (6.35,-6.25) {Ask an open-ended question};
\node[uc] (u7) at (6.35,-7.5) {Capture a crop photograph};
\node[uc] (u8) at (6.35,-8.75) {Diagnose a crop photograph};
\node[uc] (u9) at (6.35,-10.0) {Review or delete diagnoses};
\node[uc] (u10) at (6.35,-11.25) {Add notes and documents};
\node[uc] (u11) at (6.35,-12.5) {Send sensor telemetry};
\actor{0.6}{-5.5}{Farmer}
\actor{12.1}{-11.8}{Rover}
\actor{12.1}{-4.6}{Gemini and\\web search}
\foreach \u in {u1,u2,u3,u4,u5,u6,u7,u8,u9,u10} {\draw (0.9,-6.3) -- (\u.west);}
\draw (11.8,-13.2) -- (u11.east); \draw (11.8,-13.2) -- (u1.east);
\draw (11.8,-5.6) -- (u8.east); \draw (11.8,-5.6) -- (u5.east); \draw (11.8,-5.6) -- (u6.east);
\end{tikzpicture}}

\caption{Use-case diagram.}
\label{fig:usecase}
\end{figure}

\subsection{Data Model}

Figure~\ref{fig:er} shows the relational schema. Sensor readings are stored in \emph{long format} (one row per sensor per
reading), so a new sensor needs no schema change. Chunks are linked logically to their sources by a source type and a source
identifier, not by foreign keys, so one table can hold every kind of chunk. Vectors have 768 dimensions and are indexed with
HNSW (cosine); a GIN index on the text supports full-text search.

\begin{figure}[H]
\centering
\scaledtikz{\begin{tikzpicture}
\tikzset{ent/.style={nohyph,draw,rectangle split,rectangle split parts=2,align=left,font=\fontsize{8}{9.5}\selectfont,fill=white,inner sep=3pt,text width=47mm,rectangle split part fill={black!15,white}}}
\node[ent] (rd) at (0,0) {\textbf{sensor\_readings}\nodepart{second}id, device\_id, sensor, ts, value,\\received\_at (unique per\\device, sensor, ts)};
\node[ent] (wi) at (0,-3.0) {\textbf{sensor\_windows}\nodepart{second}device\_id, sensor, window\_start,\\window\_end, count, min, max,\\avg, std, first and last value};
\node[ent] (dw) at (0,-5.9) {\textbf{dirty\_windows}\nodepart{second}device\_id, window\_start,\\marked\_at};
\node[ent] (se) at (0,-8.0) {\textbf{sensor\_state\_events}\nodepart{second}device\_id, toggle,\\state (on or off), ts};
\node[ent] (hb) at (0,-10.1) {\textbf{rover\_heartbeats}\nodepart{second}device\_id, bucket\_start (5 min),\\last\_seen};
\node[ent] (kc) at (8.2,-3.0) {\textbf{knowledge\_chunks}\nodepart{second}id, source\_type, source\_id,\\chunk\_index, title, content,\\embedding vector(768),\\embed\_model, time\_start,\\time\_end, meta (JSON),\\content\_hash};
\node[ent] (fd) at (15.6,0) {\textbf{farm\_documents}\nodepart{second}id, kind (note or document),\\title, filename, content,\\chunk\_count};
\node[ent] (dg) at (15.6,-3.2) {\textbf{diagnoses}\nodepart{second}id, source, image\_path,\\status, stage, analysis (JSON),\\findings (JSON), sources,\\agent\_trace, timings};
\node[ent] (cs) at (15.6,-6.8) {\textbf{chat\_sessions}\nodepart{second}id, title, created and\\updated time};
\node[ent] (cm) at (15.6,-9.0) {\textbf{chat\_messages}\nodepart{second}id, session\_id (FK), role,\\mode, content, meta (JSON)};
\draw[arr] (rd) -- node[lbl,right]{aggregated into} (wi);
\draw[arr] (wi.east) -- node[lbl,above]{text, embedding} (wi.east -| kc.west);
\draw[arr] (fd.west) -- ++(-1.0,0) |- node[lbl,pos=0.25,left]{chunked} ([yshift=0.5cm]kc.east);
\draw[arr] (dg.west) -- ++(-1.0,0) |- node[lbl,pos=0.25,left]{findings} ([yshift=-0.5cm]kc.east);
\draw[arr] (cs) -- node[lbl,right]{1 to many} (cm);
\node[zone,fit=(rd)(se)(hb),inner sep=10pt] (zs) {};
\node[font=\fontsize{7.5}{9}\selectfont,anchor=west,text width=48mm,align=left] at (4.0,-8.4) {The readings, switch events and heartbeats together explain every gap in the data.};
\end{tikzpicture}}

\caption{Entity--relationship diagram of the database.}
\label{fig:er}
\end{figure}

\section{System Workflow}

\subsection{Data Lifecycle from Reading to Searchable Chunk}

The ESP32 samples the Nano's latest reading every five seconds, keeps up to about 30 minutes of readings in a ring buffer while
the backend is unreachable, and posts a batch every 15 seconds. The backend validates and stores each value immediately, so the
newest numbers are readable through the API at once. Closed windows are aggregated and turned into text, and the
embedding makes them searchable between 90 and 150 seconds after the window ends (a 90-second grace period for late data
plus a background job that runs every 60 seconds). Every chat question also forces an aggregation of the still-open window, so
\emph{exact statistics} are always current (Figure~\ref{fig:lifecycle}).

\begin{figure}[H]
\centering
\scaledtikz{\begin{tikzpicture}
\draw[arr] (0,0) -- (15.8,0);
\foreach \x in {1.2,4.6,8.0,11.4,14.6} {\draw[thick] (\x,-0.12) -- (\x,0.12);}
\node[box,text width=27mm,anchor=south] at (1.2,0.5) {\textbf{Reading}\\every 5 s on the\\ESP32, queued};
\node[box,text width=27mm,anchor=north] at (4.6,-0.5) {\textbf{Batch posted}\\every 15 s:\\POST /telemetry};
\node[box,text width=27mm,anchor=south] at (8.0,0.5) {\textbf{Raw rows stored}\\readable at once via\\/telemetry/latest};
\node[box,text width=27mm,anchor=north] at (11.4,-0.5) {\textbf{Window closes}\\15-minute window\\ends, 90 s grace};
\node[box,text width=27mm,anchor=south] at (14.6,0.5) {\textbf{Chunk searchable}\\statistics, text and\\embedding ready};
\node[font=\scriptsize,anchor=north] at (1.2,-0.2) {$t_0$};
\node[font=\scriptsize,anchor=north] at (8.0,-0.2) {$\approx t_0+15$ s};
\node[font=\scriptsize,anchor=north,align=center] at (14.6,-0.2) {window end\\$+\,90$ to $150$ s};
\node[font=\scriptsize,anchor=east] at (15.8,-2.6) {(not to scale)};
\end{tikzpicture}}

\caption{Lifecycle of a sensor reading.}
\label{fig:lifecycle}
\end{figure}

Three layers of text are produced from the numbers. All are written by deterministic templates:

\begin{enumerate}
\item a \textbf{window chunk} for each 15-minute window (statistics per sensor and the interpretation band of its average);
\item a \textbf{daily summary} (extremes with their times, a time-of-day profile and threshold events);
\item a \textbf{sensor on/off log} for the day (when the rover was connected and which sensors were on or off).
\end{enumerate}

An optional fourth chunk, an AI-written interpretation of the day, can be enabled; it is labelled as interpretation and is
written only from the daily summary text.

\subsection{Pipeline A: Query My Farm}

Figure~\ref{fig:flowa} shows the steps. The question is first turned into a structured \emph{query plan} by the language model (or by a
rule-based planner if no model is available): the intent (current status, statistics, comparison, trend, diagnosis history,
farm knowledge, general advice or small talk), the sensors involved, and time \emph{tokens} such as \texttt{day:-1} or
\texttt{daypart:0:morning}. Deterministic code resolves these tokens to exact ranges in the farm's time zone, so the model never
computes a date. Exact statistics become numbered \emph{facts} (\texttt{[F1]}, \texttt{[F2]}\ldots), hybrid retrieval returns
numbered \emph{chunks} (\texttt{[K1]}\ldots), and the answer-generation call receives only this context. Citation identifiers that are
not in the context are removed from the answer, and the response carries flags that say whether the answer used farm data,
the knowledge store, or general knowledge.

\begin{figure}[H]
\centering
\scaledtikz{\begin{tikzpicture}
\node[gbox,text width=62mm] (s1) at (0,0) {1. Question and conversation history};
\node[box,text width=62mm] (s2) at (0,-2.0) {2. Ingestion sync, including the\\still-filling 15-minute window};
\node[box,text width=62mm] (s3) at (0,-4.1) {3. Query plan: intent, sensors, time tokens,\\standalone rewrite (LLM, rule-based fallback)};
\node[box,text width=62mm] (s4) at (0,-6.2) {4. Resolve time tokens to farm-time ranges\\(deterministic code)};
\node[box,text width=62mm] (s5) at (0,-8.3) {5. Exact SQL facts [F1], [F2], ... and\\sensor on/off coverage facts};
\node[box,text width=62mm] (s6) at (8.4,-8.3) {6. Hybrid retrieval [K1], [K2], ...\\pgvector + full-text, fused by RRF};
\node[box,text width=62mm] (s7) at (8.4,-6.2) {7. Optional multi-hop: sensor conditions\\around a retrieved diagnosis};
\node[box,text width=62mm] (s8) at (8.4,-4.1) {8. Grounded generation: structured JSON,\\answer text streamed token by token};
\node[box,text width=62mm] (s9) at (8.4,-2.0) {9. Citation validation: remove IDs that\\are not in the context};
\node[gbox,text width=62mm] (s10) at (8.4,0) {10. Answer, sources, grounding flags;\\message saved to chat memory};
\draw[arr] (s1) -- (s2); \draw[arr] (s2) -- (s3); \draw[arr] (s3) -- (s4); \draw[arr] (s4) -- (s5);
\draw[arr] (s5) -- (s6); \draw[arr] (s6) -- (s7); \draw[arr] (s7) -- (s8); \draw[arr] (s8) -- (s9); \draw[arr] (s9) -- (s10);
\end{tikzpicture}}

\caption{Workflow of Pipeline A (``Query my farm'').}
\label{fig:flowa}
\end{figure}

\noindent\textbf{Open-ended mode.} The model plans zero to three web queries; the code runs them in parallel, removes duplicates,
ranks institutional sources (universities, government and extension services) first, labels each result's quality and drops
forum and social-media results when better sources exist. The answer cites results as \texttt{[W1]}, \texttt{[W2]}\ldots\ and the
citations are validated in the same way.

\subsection{Pipeline B: Crop Diagnosis}

Figure~\ref{fig:flowb} shows the steps. One vision-language call returns strict JSON. If the image is not a plant, the job ends.
Otherwise one research agent searches the web: healthy plants lead to cultivation and care information, unhealthy plants to
the condition, treatment and prevention, and \emph{uncertain} results are treated like unhealthy ones plus a search for how to
tell the suspected condition from its look-alikes. The agent is bounded to two rounds of plan, search and observe. A
final call writes the findings with \texttt{[W\#]} citations. The result is shown to the farmer and then written one way into the
knowledge store.

\begin{figure}[H]
\centering
\scaledtikz{\begin{tikzpicture}
\node[gbox,text width=62mm] (s1) at (0,0) {1. Photograph from the rover (/capture)\\or from the phone camera};
\node[box,text width=62mm] (s2) at (0,-2.0) {2. Validate, fix orientation, resize to\\at most 1600 px, store the file};
\node[box,text width=62mm] (s3) at (0,-4.0) {3. One VLM call: strict JSON (plant, status,\\evidence, confidence)};
\node[diamond,draw,aspect=2.6,font=\dgfont,align=center,inner sep=1pt,fill=white] (d1) at (0,-6.3) {Is it a\\plant image?};
\node[box,text width=36mm] (stop) at (-5.4,-6.3) {Stop: status\\``not\_plant''};
\node[box,text width=62mm] (s4) at (8.4,-6.3) {4. Branch policy: healthy leads to care\\information; unhealthy or uncertain\\leads to treatment information};
\node[box,text width=62mm] (s5) at (8.4,-4.0) {5. Agent loop (at most 2 rounds):\\plan queries, search, observe};
\node[box,text width=62mm] (s6) at (8.4,-2.0) {6. Write findings with [W] citations,\\remove invented citations};
\node[gbox,text width=62mm] (s7) at (8.4,0) {7. Show the result, then write the diagnosis\\and finding chunks (one-way) to the store};
\draw[arr] (s1) -- (s2); \draw[arr] (s2) -- (s3); \draw[arr] (s3) -- (d1);
\draw[arr] (d1) -- node[lbl,above]{no} (stop);
\draw[arr] (d1) -- node[lbl,above]{yes} (s4);
\draw[arr] (s4) -- (s5); \draw[arr] (s5) -- (s6); \draw[arr] (s6) -- (s7);
\end{tikzpicture}}

\caption{Workflow of Pipeline B (crop diagnosis).}
\label{fig:flowb}
\end{figure}

A diagnosis runs as a background job so that slow mobile networks never hold a long HTTP request: the upload returns at once
and the app polls the job's stage (Figure~\ref{fig:statejob}).

\begin{figure}[H]
\centering
\scaledtikz{\begin{tikzpicture}
\tikzset{st/.style={nohyph,draw,rounded corners=6pt,minimum height=8mm,minimum width=22mm,align=center,font=\dgfont,fill=white}}
\node[circle,fill=black,minimum size=3mm,inner sep=0] (init) at (0,0) {};
\node[st] (q) at (2.2,0) {queued};
\node[st] (v) at (6.4,0) {vlm};
\node[st] (s) at (10.6,0) {searching};
\node[st] (w) at (14.8,0) {writing};
\node[st,fill=black!10] (np) at (6.4,-3.3) {not\_plant};
\node[st,fill=black!10] (fail) at (10.6,-3.3) {failed};
\node[st,fill=black!10] (done) at (14.8,-3.3) {completed};
\draw[arr] (init) -- (q);
\draw[arr] (q) -- node[lbl,above]{worker starts} (v);
\draw[arr] (v) -- node[lbl,above]{plant image} (s);
\draw[arr] (s) -- node[lbl,above]{findings ready} (w);
\draw[arr] (w) -- node[lbl,right]{chunks written} (done);
\draw[arr] (v) -- node[lbl,right]{not a plant} (np);
\draw[darr] (v.south east) -- (fail.north west);
\draw[darr] (s.south) -- (fail.north);
\draw[darr] (w.south west) -- (fail.north east);
\node[font=\fontsize{8}{9.5}\selectfont,anchor=west,text width=70mm] at (0,-4.6) {Dashed arrows: any step can fail. A server restart marks every job still processing as failed, so no job stays unfinished.};
\end{tikzpicture}}

\caption{State machine of a diagnosis job.}
\label{fig:statejob}
\end{figure}

\subsection{Streaming a Chat Answer}

The chat endpoint streams Server-Sent Events~\cite{whatwgsse}: the stage names (syncing, planning, retrieving, generating), the answer text as
it is written, and finally the complete answer with validated citations, which replaces the streamed draft. The model still
returns one structured JSON object in a single call; the backend decodes the \texttt{answer} field of this JSON as it arrives, so
the phone shows words immediately without a second model call (Figure~\ref{fig:sequence}).

\begin{figure}[H]
\centering
\scaledtikz{\begin{tikzpicture}[x=2.5cm,y=-0.74cm]
\foreach \i/\n in {0/Farmer app,1/Backend API,2/Query planner,3/SQL facts,4/Knowledge store,5/Gemini} {
  \node[box,minimum width=21mm] at (\i,0) {\n};
  \draw[dashed] (\i,0.55) -- (\i,15.9);
}
\draw[arr] (0,1.3) -- node[lbl,above]{POST /chat/stream} (1,1.3);
\draw[arr] (1,2.3) -- node[lbl,above]{ingestion sync} (4,2.3);
\draw[arr] (1,3.3) -- node[lbl,above]{plan(question, history)} (2,3.3);
\draw[arr] (2,4.3) -- node[lbl,above]{structured request} (5,4.3);
\draw[darr] (5,5.3) -- node[lbl,above]{intent, sensors, time tokens} (2,5.3);
\draw[darr] (2,6.3) -- node[lbl,above]{plan} (1,6.3);
\draw[arr] (1,7.3) -- node[lbl,above]{exact statistics, coverage} (3,7.3);
\draw[arr] (1,8.3) -- node[lbl,above]{hybrid search (vector + full-text, RRF)} (4,8.3);
\draw[darr] (4,9.3) -- node[lbl,above]{[K] chunks} (1,9.3);
\draw[arr] (1,10.3) -- node[lbl,above]{context + question (stream)} (5,10.3);
\draw[darr] (5,11.3) -- node[lbl,above]{JSON text chunks} (1,11.3);
\draw[darr] (1,12.3) -- node[lbl,above]{event: stage, delta} (0,12.3);
\draw[arr] (1,13.1) -- (1.35,13.1) -- (1.35,13.8) -- (1,13.8);
\node[lbl,anchor=west] at (1.4,13.45) {validate citations};
\draw[darr] (1,14.6) -- node[lbl,above]{event: final (answer, sources)} (0,14.6);
\draw[arr] (1,15.5) -- node[lbl,above]{save message} (4,15.5);
\end{tikzpicture}}

\caption{Sequence diagram of a streamed ``Query my farm'' answer.}
\label{fig:sequence}
\end{figure}

\subsection{Sensor Switches and Data Gaps}

Each sensor can be switched on or off from the app, which saves battery because the MQ heaters dominate the rover's sensor load.
The rover reports its switches with every reading and every push. A push with no readings is a \emph{heartbeat}, which shows
that the rover is connected even when every sensor is off (video-only mode). From the switch events, the readings and the
heartbeats, every sensor in every five-minute bucket gets one status (Figure~\ref{fig:statesensor}): \emph{recorded},
\emph{off}, \emph{no reading} (switched on but silent: the heater is warming up or the soil probe is not lowered),
\emph{offline} (the rover is not connected) or \emph{unknown} (no switch information, for example older firmware).

\begin{figure}[H]
\centering
\scaledtikz{\begin{tikzpicture}
\tikzset{st/.style={nohyph,draw,rounded corners=6pt,minimum height=10mm,minimum width=32mm,align=center,font=\dgfont,fill=white,text width=30mm}}
\node[circle,fill=black,minimum size=3mm,inner sep=0] (init) at (0,0) {};
\node[st,fill=black!10] (unk) at (2.7,0) {unknown\\(no switch data)};
\node[st] (off) at (9.4,3.6) {off\\(switched off)};
\node[st] (nr) at (9.4,0) {no reading\\(on: warming or waiting)};
\node[st,fill=black!10] (rec) at (16.0,0) {recorded\\(readings arrive)};
\node[st,dashed] (offl) at (9.4,-3.8) {offline\\(rover not connected)};
\draw[arr] (init) -- (unk);
\draw[arr] (unk) -- node[lbl,above]{first report: on} (nr);
\draw[arr] (unk.north) |- node[lbl,pos=0.78,above]{first report: off} (off.west);
\draw[arr] (nr.east) to[bend left=18] node[lbl,above]{reading in bucket} (rec.west);
\draw[arr] (rec.west) to[bend left=18] node[lbl,below]{no reading} (nr.east);
\draw[arr] (nr.north) -- node[lbl,right]{switch off} (off.south);
\draw[arr] (off.east) -| node[lbl,pos=0.25,above]{switch on} (rec.north);
\draw[darr] (nr.south) -- node[lbl,right]{no heartbeat} (offl.north);
\draw[darr] (rec.south) |- node[lbl,pos=0.75,above]{no data, no heartbeat} (offl.east);
\node[font=\fontsize{8}{9.5}\selectfont,anchor=west,text width=62mm] at (0,-3.8) {Evaluated for every sensor and every five-minute bucket. A reading always wins: if the sensor produced data, the bucket is recorded.};
\end{tikzpicture}}

\caption{Status of a sensor in a five-minute bucket.}
\label{fig:statesensor}
\end{figure}

These statuses are used in three places. Window chunks and the daily summary state which sensors were missing and why. A
dedicated \emph{sensor on/off log} chunk lists the rover's sessions. At question time, a \emph{switches} fact states when
each requested sensor was recorded, off or silent, and attaches a tip (``switch MQ5 on in the app'') when the sensor is off. The
answer prompt instructs the model to quote these periods exactly, never to estimate values for them, and to say which period an
average covers.

\subsection{Live Activity View}

Every step publishes a short event into an in-memory log (readings received, window closed, chunk created, chunk embedded, chunk
stored, sensor switched, chat and diagnosis stages). The backend exposes this log as a snapshot and as a Server-Sent Events
stream, and the app's Live view animates the sensor-to-knowledge-store pipeline from it.

\section{Mathematical Model}

Symbols are defined in the Nomenclature section.

\subsection{Time Reconstruction}

The ESP32 has no real-time clock. A reading may carry a valid network-time stamp; otherwise it carries its uptime $u_i$, and the
batch carries the uptime $u^{\mathrm{sent}}$ at sending. The backend reconstructs
\begin{equation}
\hat t_i = t^{\mathrm{recv}} - \frac{u^{\mathrm{sent}} - u_i}{1000}\quad\text{(seconds)},\qquad u^{\mathrm{sent}}\ge u_i,
\label{eq:time}
\end{equation}
so a reading that was queued during an outage still receives the correct time. A value that is more than five minutes in the
future is rejected.

\subsection{Windows and Statistics}

Windows are aligned to the farm's local clock: $W_k=[t_0+k\Delta,\;t_0+(k+1)\Delta)$ where $t_0$ is local midnight and $\Delta$ divides
1440 minutes. For sensor $s$ with $n_{s,k}$ readings in $W_k$, the stored statistics are
\begin{equation}
\mu_{s,k}=\frac{1}{n_{s,k}}\sum_i x_{s,i},\qquad
\sigma_{s,k}=\sqrt{\frac{1}{n_{s,k}}\sum_i \left(x_{s,i}-\mu_{s,k}\right)^2},\qquad
m_{s,k}=\min_i x_{s,i},\quad M_{s,k}=\max_i x_{s,i}.
\label{eq:window}
\end{equation}
Statistics over a longer period $R$ (a day, a week) are combined without returning to the raw rows. The mean is count-weighted:
\begin{equation}
\mu_s(R)=\frac{\sum_{k\in R} n_{s,k}\,\mu_{s,k}}{\sum_{k\in R} n_{s,k}},\qquad
\min_R=\min_{k\in R} m_{s,k},\qquad \max_R=\max_{k\in R} M_{s,k}.
\label{eq:combine}
\end{equation}

\subsection{Hybrid Retrieval}

Embeddings are L2-normalised, so the cosine similarity is a dot product:
\begin{equation}
\cos(\mathbf q,\mathbf d)=\mathbf q\cdot\mathbf d,\qquad \text{pgvector distance}=1-\cos(\mathbf q,\mathbf d).
\label{eq:cosine}
\end{equation}
A vector ranking and a full-text ranking are fused with Reciprocal Rank Fusion~\cite{cormack2009rrf}:
\begin{equation}
\mathrm{RRF}(d)=\sum_{r\in\{\text{vector},\,\text{text}\}}\frac{1}{k+\mathrm{rank}_r(d)},\qquad k=60,
\label{eq:rrf}
\end{equation}
where a term is omitted if $d$ does not appear in ranking $r$. Chunks are returned in decreasing order of $\mathrm{RRF}(d)$.

\subsection{Sensor Status}

Let $\mathrm{Rec}_s$ be the set of buckets in which sensor $s$ produced readings, $\mathrm{Seen}$ the set of buckets in which the rover was
connected (a batch arrived or any reading exists), and $\mathrm{sw}_s(b)\in\{\mathrm{on},\mathrm{off},\varnothing\}$ the latest switch state at the end of bucket $b$.
\begin{equation}
\mathrm{st}_s(b)=
\begin{cases}
\text{recorded} & b\in \mathrm{Rec}_s\\
\text{offline} & b\notin \mathrm{Rec}_s,\ b\notin \mathrm{Seen}\\
\text{off} & b\in \mathrm{Seen},\ \mathrm{sw}_s(b)=\mathrm{off}\\
\text{no reading} & b\in \mathrm{Seen},\ \mathrm{sw}_s(b)=\mathrm{on}\\
\text{unknown} & \text{otherwise.}
\end{cases}
\label{eq:status}
\end{equation}
Readiness of a gas sensor and validity of a soil reading are
\begin{equation}
\mathrm{ready}_s(t)=\mathrm{on}_s(t)\ \wedge\ \bigl(t-t_{\mathrm{on},s}\ge T_w\bigr),\qquad
\mathrm{valid}_{\mathrm{soil}}(t)=\mathrm{on}_{\mathrm{soil}}\ \wedge\ \mathrm{arm}=\text{deployed}\ \wedge\ \neg\,\mathrm{moving},
\label{eq:ready}
\end{equation}
with $T_w=60$\,s and $t_{\mathrm{on},s}$ the time the sensor was last switched on.

\subsection{Calibration Metrics}

For $N$ diagnoses with self-reported confidence $p_i$ and correctness $y_i\in\{0,1\}$, grouped into bins $B_j$, the expected calibration
error~\cite{guo2017calibration} and the Brier score~\cite{brier1950verification} are
\begin{equation}
\mathrm{ECE}=\sum_{j}\frac{|B_j|}{N}\,\bigl|\,\mathrm{acc}(B_j)-\mathrm{conf}(B_j)\,\bigr|,\qquad
\mathrm{BS}=\frac{1}{N}\sum_{i=1}^{N}(p_i-y_i)^2.
\label{eq:calibration}
\end{equation}

\subsection{Rover Conversions and Energy}

The battery voltage is read through a divider and the soil probe is mapped linearly between a dry and a wet calibration value:
\begin{equation}
V_{\mathrm{bat}}=\frac{\mathrm{ADC}}{1023}\,V_{\mathrm{ref}}\,\frac{R_1+R_2}{R_2},\qquad
\mathrm{soil}\,[\%]=\mathrm{clamp}\!\left(100\,\frac{\mathrm{raw}-\mathrm{raw}_{\mathrm{dry}}}{\mathrm{raw}_{\mathrm{wet}}-\mathrm{raw}_{\mathrm{dry}}},\,0,\,100\right),
\label{eq:conversions}
\end{equation}
with $V_{\mathrm{ref}}=5$\,V, $R_1=10$\,k$\Omega$, $R_2=4.7$\,k$\Omega$, $\mathrm{raw}_{\mathrm{dry}}=620$ and $\mathrm{raw}_{\mathrm{wet}}=280$ (to be recalibrated per probe).
The ultrasonic distance in centimetres is the echo time in microseconds divided by 58, and joystick input is mixed as
$\mathrm{left}=\mathrm{throttle}+\mathrm{steer}$, $\mathrm{right}=\mathrm{throttle}-\mathrm{steer}$.

The energy saved by switching a sensor off for time $t_{\mathrm{off}}$ is $\Delta E=P_h\,t_{\mathrm{off}}$ when its supply is
switched. For the two MQ heaters, $P_h\approx 5\,\mathrm{V}\times 0.38\,\mathrm{A}=1.9$\,W. As an illustration, a pack of two parallel strings of cells with an
assumed \placeholder{2000\,mAh} per cell (replace with the datasheet value) stores about $2\times2.0\,\mathrm{Ah}\times7.4\,\mathrm{V}=29.6$\,Wh, so
keeping the gas sensors off for one hour saves about $1.9/29.6\approx 6\%$ of the pack, ignoring regulator losses.
